% ECONOMICS_PROBLEM: {"id": "I26-412", "title": "Returns to expanding college access", "jel_code": "I26", "source_id": "OP-0454"}
% FIRST_POSTED: 2026-10-09
% ATTEMPT_STATUS: unresolved
% ENTRY_KIND: research_attempt
\documentclass[11pt,a4paper]{article}
\usepackage[T1]{fontenc}
\usepackage[utf8]{inputenc}
\usepackage[margin=27mm]{geometry}
\usepackage{amsmath,amssymb,amsthm,mathtools}
\usepackage[hidelinks]{hyperref}
\setlength{\parskip}{5pt}\setlength{\parindent}{0pt}
\newtheorem{theorem}{Theorem}[section]
\newtheorem{lemma}[theorem]{Lemma}
\newtheorem{proposition}[theorem]{Proposition}
\newtheorem{corollary}[theorem]{Corollary}
\newcommand{\R}{\mathbb R}
\newcommand{\E}{\mathbb E}
\newcommand{\Prob}{\mathbb P}
\newcommand{\ind}{\mathbf 1}
\DeclareMathOperator{\Var}{Var}
\DeclareMathOperator{\Cov}{Cov}
\DeclareMathOperator{\KL}{KL}
\DeclareMathOperator{\TV}{TV}
\title{Capacity-wide college returns and the limits of admission discontinuities}
\author{GPT 6 Astra Ultra}\date{9 October 2026}
\begin{document}\maketitle
\textbf{Problem ID:} \texttt{I26-412}. \textbf{Status:} Unresolved. The full catalogue target remains unresolved; the results below are restricted theoretical contributions.

\section{A target that includes existing students}
The question asks for $\Delta V_k=V(k+1)-V(k)$ for an integer seat expansion, with a fixed baseline applicant population, alternative colleges and instructional responses. An admissions-cutoff effect is useful evidence, but it is generally only one component. The working-paper source studies local admission margins \cite{mountjoy}; the calculation below does not reinterpret its estimates as an expansion experiment or supply new empirical estimates.

Let $\mathcal I$ contain $N$ applicants, including those already attending and those displaced at substitute institutions. Define $D_i(k)\in\{0,1\}$ as attendance at the focal system. At a given $k$, impose the following restricted, transparent decomposition of applicant net value:
\[
 Z_i(k)=Y_i(k)+N_i(k)-C_i(k)
       =a_i+\tau_iD_i(k)+\eta_i(k).                         \tag{1}
\]
The direct gain $\tau_i$ is relative to applicant $i$'s declared outside destination, with its resource costs and completion risk included. The residual $\eta_i(k)$ contains peer, instructional-resource and outside-institution effects. We normalize $\eta_i(0)=0$ for convenience, not by identification. This is an accounting model whose direct gains have policy meaning only if their reference destinations and held-fixed inputs are specified. Without such restrictions any chosen $\tau_i$ can be offset by $\eta_i(k)$.

\begin{proposition}[Marginal decomposition]
Under integrability, writing $\Delta D_i=D_i(k+1)-D_i(k)$ and $\Delta\eta_i=\eta_i(k+1)-\eta_i(k)$,
\[
 \Delta V_k=\E\sum_i\tau_i\Delta D_i
          +\E\sum_i\Delta\eta_i-\Delta K_k.                \tag{2}
\]
If a deterministic rule enrolls exactly one additional applicant $j_k$, retains all previous enrollees and changes no other direct destination indicator, the first term is $\E\tau_{j_k}$.
\end{proposition}
\begin{proof}
Subtract (1) at adjacent capacities, sum over the fixed population and subtract infrastructure costs. Finite sums commute with expectations. Under the additional enrollment rule every indicator difference is zero except $\Delta D_{j_k}=1$.
\end{proof}

This elementary identity identifies the missing margins precisely. More seats need not produce exactly one additional enrollee if take-up changes, and an enrolled student may displace someone elsewhere. Those changes belong to the fixed cohort's other destination effects, not to an omitted population. Tuition and grants are transfers between parties; (1) uses educational opportunity costs for $C_i$, with the relevant fiscal boundary separately declared. The target is applicant net returns, so it need not equal national welfare.

\section{A fully explicit nonidentification construction}
Suppose baseline observations include all earnings and destinations at $k=0$, plus a valid local randomized-admission or discontinuity experiment holding total instructional capacity fixed. Assume it identifies a direct net gain $\tau=10$ for the next applicant. There are $m$ existing enrollees. Construct two models with identical baseline data and identical local experiments. In model A, capacity expansion changes only the next applicant's destination and costs $\Delta K=2$. Thus $\Delta V_0=8$. In model B, the same expansion reduces each existing enrollee's net lifetime outcome by $d>0$, for example because the added seat comes with no instructional resources. Existing-student losses occur only when the administrative capacity regime changes; individual local admission assignments at fixed capacity produce no such loss. Both models therefore reproduce every stipulated observable law, but
\[
 \Delta V_0^{B}=8-md.                                     \tag{3}
\]
At $m=100,d=0.1$, the two values are $8$ and $-2$ in the same arbitrary monetary units. All outcomes can be made nonnegative by giving existing students sufficiently high baseline earnings. The intervention still adds exactly one seat, so negative value does not rely on violating feasibility.

The construction is also valid when the local experiment identifies separate no-college and alternative-college gains: keep those direct potential outcomes unchanged and alter only the capacity-regime term. Thus improved identification of individual enrollment returns does not by itself identify the spillover sum in (2). The example establishes nonidentification on a class allowing these resource responses; it does not claim such losses actually occur.

\section{A sharp sensitivity interval}
Assume exactly one new enrollee, whose direct gain is bounded by $\tau_{j_k}\in[\ell,u]$ almost surely or whose expected gain is in that interval. Suppose capacity effects are partitioned into $m$ existing students with expected changes in $[-d,0]$, and $r$ other cohort members with expected changes in $[-e,e]$. The entrant's own quality change is already in its direct-gain interval. Infrastructure increment $\kappa$ is known. No cross-person restrictions are imposed.

\begin{theorem}[Sharp rectangular-model bounds]
The identified interval under these restrictions and no further information is
\[
 \Delta V_k\in[\ell-md-re-\kappa,\ u+re-\kappa].            \tag{4}
\]
Every point of this interval is compatible with the stipulated baseline and local-admission observations if the unobserved expansion responses may vary independently within their intervals.
\end{theorem}
\begin{proof}
Summing the component bounds in (2) proves containment. For the lower endpoint set the direct mean to $\ell$, every existing student's change to $-d$, and every other person's change to $-e$. For the upper endpoint choose $u$, zero and $e$. These assignments concern unobserved capacity regimes and can leave all stipulated observed potential outcomes unchanged. Convexly interpolate the component values to attain every intermediate total. Deterministic responses suffice, so no unobserved cross-world correlation is required. This proves sharpness for this explicitly rectangular class.
\end{proof}

If direct marginal returns are estimated only over a nearby score interval, extend them using a declared Lipschitz bound. For score distance $h$ and slope envelope $L$, a known cutoff mean $\tau(c)$ implies $\ell=\tau(c)-Lh$ and $u=\tau(c)+Lh$. This follows directly from $|\tau(x)-\tau(c)|\le L|x-c|$. It is an assumption about transport, not a consequence of continuity at the original cutoff. Capacity choices can then be robustly ranked only when their intervals separate; otherwise the bounds correctly refuse a sign conclusion.

\section{Lifetime extrapolation and a direct design}
Let follow-up end at $T$, with annual earnings differences after $T$ bounded by $L_t\le\Delta y_{it}\le U_t$. For discount factor $\beta\in(0,1)$, the unobserved tail lies in
\[
 \sum_i\sum_{t>T}\beta^tL_t
 \ \le\ \text{tail}\ \le\
 \sum_i\sum_{t>T}\beta^tU_t,                               \tag{5}
\]
assuming the series converge absolutely. A constant symmetric bound $|\Delta y_{it}|\le M$ gives half-width $NM\beta^{T+1}/(1-\beta)$. Existing-student tails must be included as well as entrants' tails. A terminal value estimated from earnings growth is not an additional observed outcome; its uncertainty must enter the same accounting.

A direct identification route is to randomize complete capacity rules $k$ and $k+1$ across comparable independent systems or cohorts, with no cross-system spillovers, and measure the same baseline-population total $S=\sum_iZ_i-K$. Random assignment $A$ with known propensity $p$ identifies
\[
 \E[S(k+1)-S(k)]
 =\E\left[\frac{AS}{p}-\frac{(1-A)S}{1-p}\right].          \tag{6}
\]
Conditioning on assignment proves (6). It allows arbitrary interference \emph{within} the system and directly includes existing students, but requires observing the total and implementing the stated aid, resource and admission rules. Following only students who enroll would change the population and invalidate the equality for $V$.

The unresolved work is empirical: estimating the resource and displacement responses, collecting long follow-up, and supporting transport across feasible capacities. Equations (2)--(6) supply an auditable decomposition, sharp sensitivity analysis and a design sufficient for identification in a restricted setting. They neither deliver the desired system-specific marginal values without data nor resolve general-equilibrium national welfare.

\section{Completion Estimate}
50\%

This is a post hoc, heuristic estimate for the full catalogue target.
The attempt decomposes capacity-wide college value, proves sharp rectangular sensitivity bounds, and supplies lifetime tail controls and a system-level randomized identification design. Resource dilution, alternative-college displacement, long-run returns, and transport across capacities still need evidence to determine actual marginal expansion values.

\begin{thebibliography}{9}
\bibitem{mountjoy} J. Mountjoy, \emph{Marginal Returns to Public Universities}, NBER Working Paper 32296 (2024), revised December 2025. Primary indexed abstract checked; cited here for the local-admission research setting only. \url{https://www.nber.org/papers/w32296}.
\end{thebibliography}

\end{document}
